% Self-contained recovery: tree-level angular sector and Higgs--vacuum reconstruction.
% Sources and scope: higgs_procedencia.md. Include after the constitutive operator.
\section{Joint reconstruction of gauge couplings, Higgs and the constitutive response}
\label{amp-higgs-seccion}

The vacuum response and the electroweak sector admit complementary readings
of the angular transport already constructed. The former retains both marked
channels; the determinant retains their common contraction; the angular Higgs
component also retains the orientation that disappears from that determinant.
Composing these readings reconstructs their spectral antecedent and expresses
impedance through the scalar self-coupling and the gauge couplings within one
and the same tree-level realization.

\subsection{Angular provenance and route signatures}

The APP--TRIT--TPK chain is preserved: the surviving prolongation determines
the route, the route produces its oriented register, and the register reader
expresses the angular signature. The coordinates $\pi=\pi_{\rm HMT}$ and
$\alpha=\alpha_{\rm HMT}$ are preceding outputs of the same generation;
$A$ and $C^*$ are their subsequent angular representations. Radians are used
in this block:
\begin{equation}
 x=\frac{\pi A_{\deg}}{180},\qquad
 y=\frac{\pi C^*_{\deg}}{180},\qquad
 A_{\deg}=1000\alpha,\qquad 0<y<x.
 \label{amp-higgs-angulos}
\end{equation}
Thus $x=50\pi\alpha/9$. The channels and sheet markings are those of
Section~\ref{iii:sec:hojas}:
\begin{equation}
 \begin{aligned}
 q_+&=e^{-x-y},& q_-&=e^{-x+y},& D&=q_+q_-=e^{-2x},\\
 P_\pm&=\frac{I\pm\mathcal R}{2},&
 T&=q_+P_++q_-P_-.
 \end{aligned}
 \label{amp-higgs-transporte}
\end{equation}
Here $\mathcal R$ is the marked self-adjoint involution, and each projector
has rank one in the basic two-sheet realization.

The duration--persistence angular reader is
\begin{equation}
 \chi_{\rm dur}(n,s)=\exp\!\left(nx+\frac{s}{6}y\right).
 \label{amp-higgs-caracter}
\end{equation}
To preserve the finite provenance of the coefficients used below, we assemble
the nominal fingerprints of the three routes of the length-$108$ schedule.
The cardinal-displacement columns are read as $(N,E,S,W)$; the following
pair retains the sheet and residue expressed by that route. For each
identified route, the table preserves the fingerprint and angular signature
expressed by the preceding chain.
\begin{center}
\small
\begin{tabular}{c c c c c}
\hline
Sector & Route and axis & $(N,E,S,W)$ & Sheet/residue & $(n,s)$\\
\hline
$W$ & $15/37$, E--W & $(23,44,11,30)$ & $(41,8)$ & $(94,-1)$\\
$Z$ & $20/23$, E--W & $(33,45,10,20)$ & $(57,7)$ & $(95,-1)$\\
$H$ & $24/29$, N--S & $(40,34,22,12)$ & $(47,10)$ & $(97,9)$\\
\hline
\end{tabular}
\end{center}
The three typed fingerprints preserve the route identifier, orientation
and provenance register of the signatures used. Angular evaluation follows
extraction from the enriched register; the table assembles its outputs to
compose the characters in this realization. The markings $W_\pm$ below
denote winding coordinates, distinct from the species name $W$:
\begin{equation}
 \begin{aligned}
 W_\pm&=6n\pm s, &
 n&=\frac{W_++W_-}{12},& s&=\frac{W_+-W_-}{2},\\
 (n,s)_W&=(94,-1), & (W_+,W_-)_W&=(563,565),\\
 (n,s)_Z&=(95,-1), & (W_+,W_-)_Z&=(569,571),\\
 (n,s)_H&=(97,9), & (W_+,W_-)_H&=(591,573).
 \end{aligned}
 \label{amp-higgs-firmas}
\end{equation}
The matrix of the first map is
$\bigl(\begin{smallmatrix}6&1\\6&-1\end{smallmatrix}\bigr)$, with determinant
$-12$ and image $\{(u,v)\in\mathbb Z^2:u+v\in12\mathbb Z\}$.
Indeed, the inverse formulas above are integral precisely on this image.
The character retains both orientations through
\begin{equation}
 \chi_{\rm dur}(n,s)
 =\left(q_+^{-W_+}q_-^{-W_-}\right)^{1/12}.
 \label{amp-higgs-caracter-canales}
\end{equation}
Taking logarithms proves the identity: the numerator is
$W_+(x+y)+W_-(x-y)=12nx+2sy$.
This identifies the twelfth root and the divisor $6$ of the reader.

\begin{proposition}[Angular ratios determined by the signatures]
On a common scale section, with the shared refinement factors that permit
their cancellation in the ratios,
\begin{equation}
 \frac{m_W^{\angle}}{m_Z^{\angle}}=e^{-x},\qquad
 \frac{m_H^{\angle}}{m_W^{\angle}}
 =e^{3x+5y/3},\qquad
 \left(\frac{m_H^{\angle}}{m_W^{\angle}}\right)^2
 =e^{6x+10y/3}.
 \label{amp-higgs-cocientes}
\end{equation}
\end{proposition}
\begin{proof}
Dividing the $W$ and $Z$ characters subtracts their signatures:
$(94,-1)-(95,-1)=(-1,0)$. For $H/W$, the difference is
$(97,9)-(94,-1)=(3,10)$. Equation~\eqref{amp-higgs-caracter} gives
$3x+(10/6)y$; squaring multiplies this exponent by two. The common scale
and shared refinements cancel before these operations.
\end{proof}
These ratios retain the identifier of the angular chart. Application to
masses with distinct sectorial refinements retains their additional factors;
the cancellation condition is part of the statement.

\subsection{Rationalized tree-level realization}

The electroweak reading of this chart defines
$c_W^2=D$, $s_W^2=1-D$ and $e_g^2=4\pi\alpha$. The positive root fixes
$g=e_g/s_W$ and $g'=e_g/c_W$. Hence
\begin{equation}
 g^2=\frac{4\pi\alpha}{1-D},\qquad
 g'^2=\frac{4\pi\alpha}{D}.
 \label{amp-higgs-calibre}
\end{equation}
The quantity $e_g$ is a dimensionless coupling in the rationalized
normalization, distinct from a dimensional charge section.
The scalar-potential convention is
\begin{equation}
 \begin{gathered}
 V(H)=-\mu_H^2H^\dagger H+\lambda_H(H^\dagger H)^2,\\
 m_H^2=2\lambda_Hv^2,\qquad m_W^2=\frac{g^2v^2}{4}.
 \end{gathered}
 \label{amp-higgs-potencial}
\end{equation}
The same $v$ occurs in both relations. Division cancels its scale and yields
$m_H^2/m_W^2=8\lambda_H/g^2$. Together with
\eqref{amp-higgs-cocientes}, this gives
\begin{equation}
 \boxed{\lambda_H=\frac{g^2}{8}\exp\!\left(6x+\frac{10}{3}y\right).}
 \label{amp-higgs-lambda}
\end{equation}
The factor $8$ arises from the two quadratic normalizations in
\eqref{amp-higgs-potencial}; the coefficients $6$ and $10/3$ arise from
the signatures. This distinction preserves the operation producing each
factor. The composition receives the preceding HMT outputs and uses this
tree-level physical realization at a common scale; a change of field
normalization, scale or scheme is also transported in its mass and coupling
relations.

\subsection{Joint inverse and exact domain}

\begin{theorem}[Angular reconstruction from gauge couplings and Higgs]
\label{amp-higgs-inversa-teorema}
For fixed $\pi>0$, the map defined by
\eqref{amp-higgs-transporte}, \eqref{amp-higgs-calibre} and
\eqref{amp-higgs-lambda}, on $\alpha>0$, $x>y>0$, is injective and has
the inverse
\begin{equation}
 \begin{aligned}
 D&=\frac{g^2}{g^2+g'^2},&
 x&=-\frac12\log D,&
 \alpha&=\frac{g^2g'^2}{4\pi(g^2+g'^2)},\\
 y&=\frac3{10}\left[\log\!\left(\frac{8\lambda_H}{g^2}\right)-6x\right].
 \end{aligned}
 \label{amp-higgs-inversa}
\end{equation}
Its exact image consists of the triples $g,g',\lambda_H>0$ such that,
with $D$ and $x$ reconstructed above,
\begin{equation}
 \frac{g^2}{8}e^{6x}<\lambda_H<\frac{g^2}{8}e^{28x/3}.
 \label{amp-higgs-dominio}
\end{equation}
\end{theorem}
\begin{proof}
The two gauge equations give
$g^2/(g^2+g'^2)=D$ and
$g^2g'^2/(g^2+g'^2)=4\pi\alpha$. The real logarithm recovers $x$,
and the scalar equation recovers $y$. The conditions $g,g'>0$ imply
$0<D<1$ and $x>0$. Under the strictly increasing logarithm,
inequality~\eqref{amp-higgs-dominio} is equivalent to
$0<\log(8\lambda_H/g^2)-6x<10x/3$, namely $0<y<x$.
Conversely, the reconstructed expressions give
$4\pi\alpha/(1-D)=g^2$, $4\pi\alpha/D=g'^2$ and
$g^2e^{6x+10y/3}/8=\lambda_H$. Positive roots recover the original
couplings. This proves necessity, sufficiency and uniqueness.
\end{proof}

The preceding inversion describes this map between readers exactly.
The section generated by HMT also preserves the relation in
\eqref{amp-higgs-angulos}, which becomes
\begin{equation}
 -\frac12\log\!\left(\frac{g^2}{g^2+g'^2}\right)
 =\frac{25}{18}\frac{g^2g'^2}{g^2+g'^2}.
 \label{amp-higgs-compatibilidad-alpha}
\end{equation}
Substituting the inverse expression for $\alpha$ in $x=50\pi\alpha/9$
proves the identity. The counter-angle likewise retains its preceding return
equation. Thus the reader image and the section produced by the generator
are distinct objects: the preceding compatibilities select the section
inside the reconstructible domain. Inversion acts after generation and
checks the outputs that have been retained.

\subsection{Higgs--impedance reconstruction}

To express the oriented content while preserving the channels, define
\begin{equation}
 \rho_{\rm or}=\frac{q_-}{q_+}=e^{2y},\qquad
 \Xi_H=\frac{8D^3\lambda_H}{g^2}.
 \label{amp-higgs-Xi}
\end{equation}
Equation~\eqref{amp-higgs-lambda} and $D^3=e^{-6x}$ give
\begin{equation}
 \begin{aligned}
 \Xi_H&=e^{10y/3}=\rho_{\rm or}^{5/3},&
 \rho_{\rm or}&=\Xi_H^{3/5},\\
 q_-&=\sqrt D\,\Xi_H^{3/10},&
 q_+&=\sqrt D\,\Xi_H^{-3/10}.
 \end{aligned}
 \label{amp-higgs-canales}
\end{equation}
The powers are the positive real roots. The exact domain of the oriented
chamber is
\begin{equation}
 0<D<1,\qquad 1<\Xi_H<D^{-5/3}.
 \label{amp-higgs-dominio-Xi}
\end{equation}
Indeed, $0<y<x$ is equivalent to $1<e^{10y/3}<e^{10x/3}$.
The lower endpoint gives $q_+=q_-$; the upper endpoint gives $q_-=1$.
For the two orientations together, $|y|<x$ is equivalent to
$D^{5/3}<\Xi_H<D^{-5/3}$.

Let $F(q)=f(q^{30})=(1-q^{90})/(1-q^{120})$ be the response already
proved in~\eqref{iii:eq:rchannels}--\eqref{iii:eq:monotone}.
\begin{corollary}[Confluence of scalar and constitutive readings]
\label{amp-higgs-impedancia-corolario}
On the domain~\eqref{amp-higgs-dominio-Xi},
\begin{equation}
 \boxed{\widehat Z=
 \frac{F(\sqrt D\,\Xi_H^{3/10})}
 {F(\sqrt D\,\Xi_H^{-3/10})}.}
 \label{amp-higgs-impedancia}
\end{equation}
The complete responses are
\begin{equation}
 \begin{aligned}
 r_-&=F(\sqrt D\,\Xi_H^{3/10}),&
 r_+&=F(\sqrt D\,\Xi_H^{-3/10}),\\
 \widehat\mu&=r_-^2,& \widehat\varepsilon&=r_+^2,&
 \widehat c&=(r_-r_+)^{-1}.
 \end{aligned}
 \label{amp-higgs-vacio}
\end{equation}
The two reconstructions also satisfy
\begin{equation}
 \frac{F^{-1}(\sqrt{\widehat\mu})}
 {F^{-1}(\sqrt{\widehat\varepsilon})}
 =\rho_{\rm or}=\Xi_H^{3/5}.
 \label{amp-higgs-confluencia}
\end{equation}
\end{corollary}
\begin{proof}
Substituting the channels of~\eqref{amp-higgs-canales} into the four
readings of~\eqref{iii:eq:four} proves the first identities. The
constitutive inverse of~\eqref{iii:eq:cubic} recovers $q_-$ and $q_+$;
their ratio is $\rho_{\rm or}$. Its equality to $\Xi_H^{3/5}$ was proved
in~\eqref{amp-higgs-canales}.
\end{proof}

On the section compatible with $A_{\deg}=1000\alpha$, the coupling
$Q(D):=-9\log D/25$ satisfies $Q(D)=4\pi\alpha$ and
$g^2=Q(D)/(1-D)$. Thus the same Higgs coordinate takes the form
\begin{equation}
 \Xi_H=\frac{8(1-D)D^3\lambda_H}{Q(D)},\qquad
 \frac1{g^2}+\frac1{g'^2}=\frac1{Q(D)}.
 \label{amp-higgs-curva}
\end{equation}
These expressions relate the total coupling, its rationalized partition
and the scalar orientation on the same section.

\subsection{Retained information and transport of orientation}

The gauge pair recovers $\alpha$ and $x$; the ratio $\lambda_H/g^2$
then recovers $y$. Together with the markings $P_\pm$, it reconstructs
the complete operators $T$ and $\mathcal V=r_+P_++r_-P_-$.
The three scalars determine the eigenvalues; the marked projectors also
determine their operator realization. The enriched route history remains
preceding information of higher resolution.

Under $y\mapsto-y$ relative to the fixed markings, $D,g,g'$ remain
invariant and $\Xi_H\mapsto\Xi_H^{-1}$. The constitutive readings are
exchanged, $\widehat Z\mapsto\widehat Z^{-1}$, and $\widehat c$ is preserved.
If $\lambda_H^+$ and $\lambda_H^-$ denote the evaluations in the two
conjugate orientations, then
\begin{equation}
 \lambda_H^+\lambda_H^-
 =\frac{g^4}{64}D^{-6}.
 \label{amp-higgs-orientacion}
\end{equation}
Multiplying~\eqref{amp-higgs-lambda} at $y$ and $-y$ proves the identity:
the oriented contributions cancel. This comparison retains the signature
reader and changes its angular argument. A simultaneous conjugation of
routes and reference markings must also transport the signatures.

For the amplification $T_m=T\otimes I_{9^m}$, with
$\tau_m=\operatorname{Tr}/(2\cdot9^m)$ and
$\mathcal R_m=\mathcal R\otimes I_{9^m}$, the appropriate coordinates are
\begin{equation}
 D=\exp\bigl(2\tau_m(\log T_m)\bigr),\qquad
 \rho_{\rm or}=\exp\bigl(-2\tau_m(\mathcal R_m\log T_m)\bigr).
 \label{amp-higgs-refinamiento}
\end{equation}
Indeed, $\log T_m=(\log T)\otimes I_{9^m}$, so the two traces are $-x$
and $-y$, respectively. Normalization removes the amplification multiplicity.
The gauge, Higgs and vacuum readings are therefore compatible with these
refinements; the ordinary amplified determinant would be $D^{9^m}$ and
has a different role.

The composition has a precise conclusion: the even gauge invariant and
the oriented scalar coordinate jointly reconstruct the spectral pair
governing the vacuum. When $\lambda_H$ is evaluated using the same formula,
the confluence is a consistency identity. An independent reading of
$\lambda_H$, transported to identical scheme and scale, makes the equality
an additional test. Both uses retain the signatures, sheet markings and
tree-level normalization explicitly.
